% ============================================================
% CAPÍTULO 5: CONCLUSIONES
% ============================================================

\chapter{Conclusiones y Líneas Futuras}
\label{chap:conclusiones}

A continuación se presental las conclusiones finales del proyecto, cumplimiento de objetivos, valoraciones
y líneas futuras de investigación derivadas de este trabajo.

\section{Resumen y aportaciones del proyecto}
\label{sec:resumen_conclusiones}

Este trabajo ha abordado el problema de la predicción inminente de la
Fibrilación Auricular Paroxística (PAF) a partir de señales de electrocardiograma en ritmo sinusal.
Las principales contribuciones de este proyecto incluyen:

\begin{itemize}
  \item \textbf{Pipeline de Preprocesamiento Armonizado y Dinámico:}
  Se ha diseñado un pipeline capaz de unificar cinco bases de datos de distinta procedencia
  (\texttt{afpdb}, \texttt{cpsc2021}, \texttt{ltafdb}, \texttt{shdb-af} y \texttt{afdb})
  gestionando discrepancias entre ellas como frecuencias de muestreo distintas, que obliga a un
  resampleo uniforme a $128$~Hz.
  \item \textbf{Evaluación Comparativa de Arquitecturas:}
  Se han comparado bajo validación cruzada por grupos (inter-sujeto)
  redes neuronales residuales (ResNet 1D), mecanismos de atención por canal (SEResNet 1D),
  autoatención temporal híbrida (CNN-Transformer) y el modelo de espacio de estados selectivo
  (Mamba 1D).
  \item \textbf{Demostración de la Sinergia Híbrida (ECG + HRV):}
  Se ha corroborado empíricamente que la inclusión conjunta de información morfológica a corto plazo
  (ECG con una ventana de $10\text{ s}$) y a mediano plazo (HRV en ventana de
  $5\text{ min}$) proporciona una mejora en la clasificación.
\end{itemize}

\section{Cumplimiento de objetivos}
\label{sec:cumplimiento_objetivos}

Se ha alcanzado con éxito el propósito principal del trabajo a través de los siguientes hitos:
\begin{itemize}
  \item \textbf{Detección del estado crítico Pre-PAF:} Se ha logrado identificar las variaciones morfológicas
  que llevan a un episodio de fibrilación auricular en la ventana crítica de 5 minutos inmediatamente anterior al episodio, logrando un
  F1-score del $0.81$ en la clase Pre-PAF mediante la configuración ResNet 1D, o si atendemos a modelos clásicos
  con un F1-score de $0.87$ con \textit{Random Forest}.
  \item \textbf{Mitigación del Data Leakage:} La separación estricta de conjuntos de datos agrupando
  por paciente garantiza que las métricas de rendimiento reportadas no estén
  sesgadas por memorización de morfologías específicas de un sujeto.
  \item \textbf{Estudio de la contribución HRV:} Se evaluó el peso relativo de la morfología del ECG y las carácteristicas HRV.
\end{itemize}

\section{Valoración crítica y limitaciones}
\label{sec:valoracion_critica}

El estudio presenta limitaciones que deben ser consideradas:
\begin{itemize}
  \item \textbf{Volumen acotado de datos clínicos:}
  El aprendizaje profundo requiere volúmenes masivos de datos para obtener los mejores resultados.
  En este proyecto, la cantidad limitada de señales
  explica que arquitecturas muy sofisticadas (como el Transformer o SEResNet) muestren un sobreajuste
  local y no consigan superar a clasificadores clásicos como el
  \textit{Random Forest} alimentado con datos HRV precalculados.
  \item \textbf{Inestabilidad frecuencial de HRV en ventanas cortas:}
  Para calcular de forma fiable la banda espectral de baja frecuencia (LF, $0.04-0.15\text{ Hz}$),
  se deberían registrar en la señal al menos dos ciclos de la oscilación más baja ($25\text{ s}$
  por ciclo), exigiendo registros mínimos de $50$ a $75\text{ s}$.
  El uso de ventanas cortas (menores a 1 minuto)
  induce una inestabilidad matemática en los descriptores espectrales debido a la pérdida
  de resolución frecuencial.
\end{itemize}

\section{Líneas futuras}
\label{sec:lineas_futuras}

A continuación, se proponen las siguientes líneas de desarrollo futuro:
\begin{itemize}
  \item \textbf{Exploración de Modelos Fundacionales Preentrenados (Signal Transformers):}
  Investigar la transferencia de aprendizaje utilizando modelos auto-supervisados entrenados
  en millones de horas de ECG clínicos, lo que resolvería la limitación del tamaño de datos local.
  \item \textbf{Validación en Tiempo Real y Wearables:} Evaluar los clasificadores embebidos en
  dispositivos en tiempo real, validando su sensibilidad ante el ruido de movimiento cotidiano.
\end{itemize}
